%\documentclass{article}
\documentclass{exam}
\author{}


\date{}
\usepackage{amsmath}
\usepackage[margin=2.5cm]{geometry}

\title{Microeconomics Cheat-Sheet}

\begin{document}

\maketitle


\section{Preliminaries}

Basic economic terminology:

% \begin{equation}
%     Y = C + I + G + NX
% \end{equation}

% where:
\begin{itemize}
    \item Real v Nominal Prices
    \item Positive v Normative Analysis
    \item Definition of a market; market segmentation
\end{itemize}

% GDP can also be viewed as total \underline{production} (hence its name) or total \underline{income} over the period.

\section{Supply \& Demand}

% The Cobb-Douglas production function is a simple model of the relation between
% inputs and output in an economy:

% \begin{equation}
%     Y = A K^{\alpha} L^{1-\alpha}
% \end{equation}

% where:

Some more terminology:

\begin{itemize}
    \item Supply Curve, Demand Curve
    \item Market mechanism and equilibrium price and quantity
    \item Elasticity in general
    \item Substitute goods and complementary goods
\end{itemize}

The Cobb-Douglas formula is perhaps the simplest possible function of K and L
which satisfies three reasonable requirements:

\begin{itemize}
    \item Increases in inputs ($K$ and/or $L$) generate increases in outputs ($Y$)
    \item There are diminishing returns to $K$ and $L$
    \item There are constant returns to scale - if $K$ and $L$ both rise by a factor of
          $\lambda$ then $Y$ rises by $\lambda$ too.
\end{itemize}


Frequently it's more useful to express Y and K on a per-labor-unit ($L$) basis.  $Y/L$ is called
the "average product of labor" and $K/L$ is called "capital intensity."
Dividing both sides of equation (2) by $L$ produces

\begin{equation}
    \frac{Y}{L} = A \left(\frac{K}{L}\right)^{\alpha}
\end{equation}

If we write $Y/L$ as $y$ and $K/L$ as $k$, the formula becomes very simple:

\begin{equation}
    y = A k^{\alpha}
\end{equation}

\section{The Solow Growth Model}

In the Solow model, capital accumulation is determined by the savings rate and
the depreciation rate. Its simplifying assumption is that savings equals investment, and that
savings is a constant percentage of output ($s$).  The equation for capital accumulation from
time $t$ to time $t+1$ is:

\begin{equation}
    K_{t+1} = (1-\delta) K_t + sY_t
\end{equation}

where:
\begin{itemize}
    \item $K$ is the capital stock,
    \item $\delta$ is the depreciation rate of capital,
    \item $s$ is the savings rate (the portion of production converted to capital),
    \item $Y$ is GDP
\end{itemize}

By substituting $Y = A K^{\alpha} L^{1-\alpha}$ for $Y$ in equation 5, rearranging terms, and setting
$K_{t+1} = K_t$, we get a number of useful "steady-state" identities:

\begin{itemize}
    \item Capital per worker:  $\frac{K_{ss}}{L}=\left(\frac{s}{\delta}A\right)^\frac{1}{1-\alpha}$
    \item Capital to output: $\left(\frac{K}{L}\right)_{ss} = \frac{s}{\delta}$
    \item Growth rate in $\frac{Y}{L}$:  $\frac{(1+g_Y)}{(1+g_L)} = (1+g_A)^\frac{1}{1-\alpha}$
\end{itemize}


That last identify makes it clear that the steady-state growth rate relies completely on $A$ and
is independent of $s$ and $\delta$.

\section{Unemployment}

The progression of unemployment $u$ can be modeled as

\begin{equation}
    u_{t+1} = u_t + se_t - au_t
\end{equation}

where:
\begin{itemize}
    \item $u$ is the unemployment rate
    \item $e$ is the employment rate (ignoring inactives)
    \item $s$ is the separation rate (probability of termination)
    \item $a$ is the accession rate (probability of hire)
\end{itemize}

Substituting $(1-u)$ for $e$ and solving for the steady-state level of
unemployment, we get:

\begin{equation}
    u_{ss} = \frac{s}{s+a}
\end{equation}


\end{document}
